\section{Elektrochemie}

Elektrochemie je obor chemie, který zkoumá procesy probíhající na rozhraní elektrod a elektrolytu.

\subsection{Beketovova řada kovů}

Pojmenována po ruském chemikovi N. N. Beketovovy. Kovy jsou seřazeny podle hodnoty standardního elektrodového potenciálu.

\begin{tabular}{|*{9}{c|}}
	\hline
	Li/\ce{Li+} & K/\ce{K+} & Mg/\ce{Mg^{2+}} & Zn/\ce{Zn^{2+}} & Cd/\ce{Cd^{2+}} & \textbf{H/}\ce{\textbf{\ce{H+}}} & Cu/\ce{Cu^{+}} & Hg/\ce{Hg^{2+}} & Au/\ce{Au^{3+}} \\
	\hline
	$-$3,04 & $-$2,93 & $-$2,37 & $-$0,761 & $-$0,40 & \textbf{0} & +0,52 & +0,85 & +1,52 \\
	\hline
\end{tabular}

Kovy, které mají nižší potenciál než vodík se označují jako \textit{neušlechtilé}, kovy s vyšším potenciálem jsou pak označovány jako \textit{ušlechtilé}.

Neušlechtilé kovy při reakci s minerálními kyselinami uvolňují vodík:

\ce{Zn + H2SO4 -> ZnSO4 + H2}

Ušlechtilé kovy nedokáží vodík vytěsnit, rozpouštějí se jen v oxidačních kyselinách:

\ce{3 Cu + 8 HNO3 -> 3 Cu(NO3)2 + 2 NO + 4 H2O}

\textbf{Standardní vodíková elektroda} -- platinová elektroda, pokrytá platinovou černí, sycená plynným vodíkem pod tlakem 101,325 kPa za teploty 0 $^\circ$C, ponořená do roztoku s jednotkovou aktivitou vodíkových iontů.

\ce{\nicefrac{1}{2} H2 <=> H+ + e-}

\newpage

\subsection{Elektrody a elektrodový potenciál}

$E = E^0 - \frac{RT}{zF}\ln a$

$E = E^0 - \frac{RT}{zF}\ln \frac{a_{red}}{a_{ox}}$

\subsection{Galvanický článek}

\subsection{Elektrolýza}

\textbf{1. Faradayův zákon}\\
Hmotnost látky vyloučené na elektrodě závisí přímo úměrně na elektrickém proudu, procházejícím elektrolytem, a na čase, po který elektrický proud procházel.

m = A . I . t

\textbf{2. Faradayův zákon}\\
Látková množství vyloučená stejným nábojem jsou pro všechny látky chemicky ekvivalentní, neboli elektrochemický ekvivalent A závisí přímo úměrně na molární hmotnosti látky.

$A = \frac{M}{F . z}$